% GENERATED FILE. Edit canonical lessons, then run npm run generate:books.
% canonical-source-hash: fnv1a64:51d424a51cded0d3
% canonical-lessons: GE-C350-arbeitgeber, GE-C350-arbeitnehmer, GE-C350-einstellen, GE-C350-beschaeftigen, GE-C350-befristet

\chapter{Arbeitgeber und Arbeitnehmer}
\label{ch:ge-arbeitgeber-und-arbeitnehmer}
\hlchaptermodality{350}

\begin{chapteropening}
\textbf{By the end of this chapter:} \emph{I can talk about employers and employees, hiring, and contracts with an end date.}

Complete the last lesson of chapter 350: I can talk about employers and employees, hiring, and contracts with an end date.
\end{chapteropening}

\section[der Arbeitgeber]{der Arbeitgeber — the employer}

\label{lesson:GE-C350-arbeitgeber}



\begin{quote}
Before the new one: say the German for qualified, then the German for to master.
\end{quote}

\subsection*{You'll want to know: der Arbeitgeber}
\textbf{der Arbeitgeber} — \textquotedblleft{}the employer\textquotedblright{}.

\begin{quote}
Mein Arbeitgeber bezahlt mir das Ticket für den Bus.
\end{quote}

\begin{quote}
\emph{My employer pays for my bus ticket.}
\end{quote}

Plural: \emph{die Arbeitgeber}; a woman is \emph{die Arbeitgeberin}. Literally the work-giver.

\begin{cousinweb}
\emph{Arbeitgeber} joins \emph{Arbeit} \textquotedblleft{}work\textquotedblright{}, from Old High German \emph{arbeit} and Proto-Germanic \emph{arbaidiz} \textquotedblleft{}labour, hardship\textquotedblright{}, to \emph{Geber} \textquotedblleft{}giver\textquotedblright{}, from \emph{geben} \textquotedblleft{}to give\textquotedblright{} (EtymDB).
\end{cousinweb}

\begin{grammarlens}[title={What you've built}]
One more word for telling, explaining and describing at B1.
\end{grammarlens}

\subsection*{Your turn}
\begin{itemize}
\raggedright
  \item \emph{Say it:} \emph{der Arbeitgeber}
  \item \emph{Say it:} \emph{der Arbeitgeber}, once more
  \item \emph{Say it:} a sentence of your own with \emph{der Arbeitgeber}
  \item \emph{Recall:} say the German for qualified, then the German for to master, then say \emph{der Arbeitgeber} again
\end{itemize}

\begin{tcolorbox}[breakable,colback=teal!4,colframe=teal!35!black,title={Before you move on}]
What does der Arbeitgeber mean? (\textquotedblleft{}The employer\textquotedblright{}.) Say it once more.
\end{tcolorbox}
\section[der Arbeitnehmer]{der Arbeitnehmer — the employee}

\label{lesson:GE-C350-arbeitnehmer}



\begin{quote}
Before the new one: say the German for to master, then the German for the employer.
\end{quote}

\subsection*{You'll want to know: der Arbeitnehmer}
\textbf{der Arbeitnehmer} — \textquotedblleft{}the employee\textquotedblright{}.

\begin{quote}
Die Firma hat zweihundert Arbeitnehmer.
\end{quote}

\begin{quote}
\emph{The company has two hundred employees.}
\end{quote}

Plural: \emph{die Arbeitnehmer}; a woman is \emph{die Arbeitnehmerin}. Literally the work-taker.

\begin{cousinweb}
\emph{Arbeitnehmer} joins \emph{Arbeit} \textquotedblleft{}work\textquotedblright{}, Proto-Germanic \emph{arbaidiz} (EtymDB), to \emph{Nehmer} \textquotedblleft{}taker\textquotedblright{}, from \emph{nehmen} \textquotedblleft{}to take\textquotedblright{}, Old High German \emph{neman} (EtymDB).
\end{cousinweb}

\begin{grammarlens}[title={What you've built}]
One more word for telling, explaining and describing at B1.
\end{grammarlens}

\subsection*{Your turn}
\begin{itemize}
\raggedright
  \item \emph{Say it:} \emph{der Arbeitnehmer}
  \item \emph{Say it:} \emph{der Arbeitnehmer}, once more
  \item \emph{Say it:} a sentence of your own with \emph{der Arbeitnehmer}
  \item \emph{Recall:} say the German for to master, then the German for the employer, then say \emph{der Arbeitnehmer} again
\end{itemize}

\begin{tcolorbox}[breakable,colback=teal!4,colframe=teal!35!black,title={Before you move on}]
What does der Arbeitnehmer mean? (\textquotedblleft{}The employee\textquotedblright{}.) Say it once more.
\end{tcolorbox}
\section[einstellen]{einstellen — to hire, to take on; to adjust; to stop}

\label{lesson:GE-C350-einstellen}



\begin{quote}
Before the new one: say the German for the employer, then the German for the employee.
\end{quote}

\subsection*{You'll want to know: einstellen}
\textbf{einstellen} — \textquotedblleft{}to hire, to take on; to adjust; to stop\textquotedblright{}.

\begin{quote}
Das Restaurant stellt zwei neue Köche ein.
\end{quote}

\begin{quote}
\emph{The restaurant is hiring two new cooks.}
\end{quote}

Separable: \emph{sie stellt ein}, \emph{sie hat eingestellt}. Also to adjust: \emph{die Uhr einstellen}; and to stop: \emph{die Arbeit einstellen}.

\begin{cousinweb}
\emph{einstellen} is built from \emph{ein} \textquotedblleft{}in\textquotedblright{} and \emph{stellen} \textquotedblleft{}to put, to place\textquotedblright{}, inherited from Old High German \emph{stellen} and Proto-Germanic \emph{stallijaną} (EtymDB): you put someone into a job.
\end{cousinweb}

\begin{grammarlens}[title={What you've built}]
One more word for telling, explaining and describing at B1.
\end{grammarlens}

\subsection*{Your turn}
\begin{itemize}
\raggedright
  \item \emph{Say it:} \emph{einstellen}
  \item \emph{Say it:} \emph{einstellen}, once more
  \item \emph{Say it:} a sentence of your own with \emph{einstellen}
  \item \emph{Recall:} say the German for the employer, then the German for the employee, then say \emph{einstellen} again
\end{itemize}

\begin{tcolorbox}[breakable,colback=teal!4,colframe=teal!35!black,title={Before you move on}]
What does einstellen mean? (\textquotedblleft{}To hire, to take on; to adjust; to stop\textquotedblright{}.) Say it once more.
\end{tcolorbox}
\section[beschäftigen]{beschäftigen — to employ; to occupy (sich beschäftigen mit)}

\label{lesson:GE-C350-beschaeftigen}



\begin{quote}
Before the new one: say the German for the employee, then the German for to hire.
\end{quote}

\subsection*{You'll want to know: beschäftigen}
\textbf{beschäftigen} — \textquotedblleft{}to employ; to occupy (sich beschäftigen mit)\textquotedblright{}.

\begin{quote}
Die kleine Bäckerei beschäftigt zehn Leute.
\end{quote}

\begin{quote}
\emph{The small bakery employs ten people.}
\end{quote}

Inseparable: \emph{sie beschäftigt, sie hat beschäftigt}. \emph{Sich mit etwas beschäftigen} — to deal with something. \emph{Beschäftigt} — busy — is the same family.

\begin{cousinweb}
\emph{beschäftigen} comes from Middle High German \emph{scheftec} \textquotedblleft{}busy, active\textquotedblright{} (EtymDB), of the family of \emph{schaffen} \textquotedblleft{}to make, to create\textquotedblright{}: you give people work to do.
\end{cousinweb}

\begin{grammarlens}[title={What you've built}]
One more word for telling, explaining and describing at B1.
\end{grammarlens}

\subsection*{Your turn}
\begin{itemize}
\raggedright
  \item \emph{Say it:} \emph{beschäftigen}
  \item \emph{Say it:} \emph{beschäftigen}, once more
  \item \emph{Say it:} a sentence of your own with \emph{beschäftigen}
  \item \emph{Recall:} say the German for the employee, then the German for to hire, then say \emph{beschäftigen} again
\end{itemize}

\begin{tcolorbox}[breakable,colback=teal!4,colframe=teal!35!black,title={Before you move on}]
What does beschäftigen mean? (\textquotedblleft{}To employ; to occupy (sich beschäftigen mit)\textquotedblright{}.) Say it once more.
\end{tcolorbox}
\section[befristet]{befristet — fixed-term, temporary}

\label{lesson:GE-C350-befristet}



\begin{quote}
Before the new one: say the German for to hire, then the German for to employ.
\end{quote}

\subsection*{You'll want to know: befristet}
\textbf{befristet} — \textquotedblleft{}fixed-term, temporary\textquotedblright{}.

\begin{quote}
Mein Vertrag ist auf ein Jahr befristet.
\end{quote}

\begin{quote}
\emph{My contract is fixed-term, for one year.}
\end{quote}

The participle of \emph{befristen} — to set a time limit. The opposite is \emph{unbefristet} — permanent.

\begin{cousinweb}
\emph{befristet} is the participle of \emph{befristen}, built on \emph{Frist} \textquotedblleft{}time limit\textquotedblright{}, inherited from Middle High German \emph{vrist}, Old High German \emph{frist} and Proto-Germanic \emph{fristiz} (EtymDB).
\end{cousinweb}

\begin{grammarlens}[title={What you've built}]
One more word for telling, explaining and describing at B1.
\end{grammarlens}

\subsection*{Your turn}
\begin{itemize}
\raggedright
  \item \emph{Say it:} \emph{befristet}
  \item \emph{Say it:} \emph{befristet}, once more
  \item \emph{Say it:} a sentence of your own with \emph{befristet}
  \item \emph{Recall:} say the German for to hire, then the German for to employ, then say \emph{befristet} again
\end{itemize}

\begin{tcolorbox}[breakable,colback=teal!4,colframe=teal!35!black,title={Before you move on}]
What does befristet mean? (\textquotedblleft{}Fixed-term, temporary\textquotedblright{}.) Say it once more.
\end{tcolorbox}
